\section{Some reverse mathematics of Knaster-Tarski}
\label{sec:rev-math-kt}
In order to obtain sub-recursive upper bounds on the functions represented by $\cmuLJ$, we will need to formalise the totality argument of the previous section itself within fragments of `second-order' arithmetic.
% 
To this end we will need to formalise some of the reverse mathematics about fixed point theorems on which our type structure relies.




\subsection{Language  and theories of `second-order' arithmetic}
Let us recall
$\langsoarith$, the language of `second-order' arithmetic, e.g.\ as given in \cite{Simpson99:monograph}.
It extends the language of arithmetic $\langarith$ by:
\begin{itemize}
    \item an additional sort of \emph{sets}, whose variables are written $X,Y $ etc. Individuals of $\langarith$ are considered of \emph{number} sort.
    \item an \emph{elementhood} (or \emph{application}) relation $\in$ relating the number sort to the set sort. I.e.\ there are formulas $t \in X$ (also $Xt$) when $t $ is a number term and $X$ is a set variable. (We will not consider any non-variable set terms here.)
\end{itemize}

When speaking about the `free variables' of a formula, we always include set variables as well as individual variables.
% 
We shall assume a De Morgan basis of connectives, namely $\lor, \land, \exists, \forall$ with negation only on atomic formulas.
Hence, we say that a formula $\phi$ is \emph{positive} (or \emph{negative}) in $X$ if no (every, resp.) subformula $Xt$ occurs under a negation.
Let us note that we have an analogue of `functoriality' in predicate logic:
\begin{lemma}[Monotonicity]
\label{lem:pos-imp-mon}
Let $\phi(X,x)$ be positive in $X$. 
% We can derive in pure logic:
\[
\proves \forall x (Xx \limp Yy) \limp \forall x (\phi(X,x) \limp \phi(Y,x))
\]
\end{lemma}
\begin{proof}[Proof sketch]
By (meta-level) induction on the structure of $\phi$.
\end{proof}
In what follows this will often facilitate arguments by allowing a form of `deep inference' reasoning.


% 
We shall work with subtheories of full second-order arithmetic ($\SOPA$) such as $ \ACA $, $  \PCA$, $  \PSCA$ etc., whose definitions may be found in standard textbooks, e.g.\ \cite{Simpson99:monograph}.
We shall freely use basic facts about them.
% from, e.g., \cite{Simpson99:monograph}.
For instance:

\begin{proposition}
    [Some basic reverse mathematics, e.g.\ \cite{Simpson99:monograph}]
    \label{prop:basic-rev-math}
    We have the following:
    \begin{enumerate}
        \item\label{item:inclusions-basic-rev-math} $ \ACA \subseteq \PCA = \SCA \subseteq \PSCA = \SPCA $
        \item\label{item:pca-proves-atr-basic-rev-math} $\PCA \proves \ATR$ (\emph{arithmetical transfinite recursion}).
        \item\label{item:psca-proves-spac-basic-rev-math} $\ac{\Din 1 2}_0 \proves \ac{\Sin 1 2}$  (\emph{axiom of choice})
    \end{enumerate}
\end{proposition}


A simple consequence of $\Sin 1 2$-choice in $\PSCA$ is that $\Sin 1 2$ (and $\Pin 1 2 $) is provably closed under positive $\Sin 1 1 $ (resp.\ $\Pin 1 1 $) combinations (even with $\Sin 1 1 $ and $\Pin 1 1 $ parameters).
We shall actually need a refinement of this fact to take account of polarity, so we better state it here.
First let us set up some notation.


\begin{definition}
    [Polarised analytical hierarchy]
    We write $\posnegSin 1 n {\vec X}{\vec Y}$ for the class of $\Sin 1 n $ formulas positive in $\vec X$ and negative in $\vec Y$.
Similarly for $\posnegPin 1 n {\vec X}{\vec Y}$.
$\phi$ is $\posnegDin 1 n {\vec X}{\vec Y}$, in a theory $T$, if it is $T$-provably equivalent to both a $\posnegSin 1 n {\vec X}{\vec Y} $-formula and a $\posnegPin 1 n{\vec X}{\vec Y}$-formula.
\end{definition}

\begin{lemma}
[Polarised substitution lemma, $\PSCA$]
\label{lem:pol-sub}
If $\phi(X) \in \posnegSin 1 1 {X,\vec X}{\vec Y}$ and $\psi\in \posnegSin 1 2 {\vec X}{\vec Y}$ then $\phi(\psi) \in \posnegSin 1 2 {\vec X}{\vec Y}$.
Similarly for $\Pi$ in place of $\Sigma$.
\end{lemma}
\begin{proof}
    [Proof sketch]
    By induction on the structure of $\phi$, using $\ac{\Sin 1 2}$ at each alternation of a FO and SO quantifier.
\end{proof}

Note that the Lemma above, in particular, allows for arbitrary substitutions of $\Sin 1 1 $ and $\Pin 1 1 $ formulas free of $\vec X,\vec Y$, which we shall rely on implicitly in the sequel.




% We shall also assume basic facts about orders and comparison within second-order arithmetic.
% We write $\WO(\alpha)$ if $\alpha$ is a (countable) well-order, $<_\alpha$ for comparison within $\alpha$, and $\ordleq$ for comparison of well-orders.


\subsection{Countable orders}
We can develop a basic theory of (countable) ordinals in even weak second-order theories, as has been done in \cite{Simpson99:monograph} and also comprehensively surveyed in \cite{hirst05:ord-arith-rm-survey}.
Let us point out that, while distinctions between natural notions of `order comparison' are pertinent for weak theories, the theories we mainly consider contain $\ATR_0$, for which order comparison is robust.

A \emph{(countable) binary relation} is a pair $ (X,\leq)$ where $X$ and $\leq$ are sets, the latter construed as ranging over pairs. 
(Sometimes we write $\Rel(X,\leq)$ to specify this.)
We say that $(X,\leq)$ is a \emph{partial order}, written $\PO(X,\leq)$, if:
\begin{itemize}
    \item $\forall x \in X \, x\leq x$
    \item $\forall x,y \in X (x\leq y \leq x \limp x=x)$
    \item $\forall x,y,z (x\leq y\leq z \limp x\leq z)$
\end{itemize}

$(X,\leq)$ is a \emph{total order}, written $\TO(X,\leq)$, if it is a partial order that is \emph{total}:
\begin{itemize}
    \item $\forall x,y (x\leq y \lor y\leq x)$
\end{itemize}

Given a relation $\leq$, we may write $<$ for its \emph{strict} version, given by
\[
x<y := x\leq y \land \lnot y\leq x
\]
We employ similar notational conventions for other similar order-theoretic binary symbols.

We say that a binary relation $(X,\leq)$ is \emph{well-founded}, written $\WF(X,\leq)$, if:
\begin{itemize}
    \item $\forall f: \Nat \to X \exists x \lnot f(x+1)<f(x)$
\end{itemize}

% It is direct from the definition that $\ACA$ admits an induction principle on well-founded relations:
% \begin{fact}
%     [Well-founded induction, $\ACA$]
%     \label{fact:wf-ind-aca}
%     Suppose $\WF(X,\leq )$. Then:
%     \[
%     \forall x (\forall y\leq x \, Xy \limp Xx) \limp \forall x Xx
%     \]
%     \anupam{consider changing inequality everywhere here to not clash with usual inequality}
% \end{fact}

$(X,\leq)$ is a \emph{well-order}, written $\WO(X,\leq)$, if it is a well-founded total order, i.e.:
\[
\WO(X,\leq) := \TO(X,\leq) \land \WF(X,\leq)
\]

Henceforth we shall write $\alpha,\beta$ etc.\ to range over countable binary relations.
If $\alpha = (X,\leq)$ we may write $x\leq_\alpha y := x,y \in X \land x\leq y$, and similarly $x<_\alpha y$ for $x,y\in X \land x<y$.
We may also write simply $x \in \alpha$ instead of $x \in X$, as abuse of notation.

It is not hard to see that $\ACA$ admits an induction principle over any provable well-order (see, e.g., \cite{Simpson99:monograph}):
\begin{fact}
[Transfinite induction, $\ACA$]
\label{fact:wo-ind-in-aca}
If $\WO(\alpha)$ then:
\[
% \left( 
\forall X (\forall x\in \alpha (\forall y<_\alpha x\,  X(y) \limp X(x)) \limp \forall x \in \alpha\,  X(x))
% \right)
\]
\end{fact}
In fact, in extensions of $\ACA$, we even have (transfinite) induction on arithmetical formulas, a fact that we shall use implicitly in what follows.
For instance, in $\PSCA$, we may admit (transfinite) induction on arithmetical combinations of $\Pin 1 2 $ formulas.
\anupam{check this again}

\subsection{Comparing orders}
Following Simpson in \cite{Simpson99:monograph},
given $\alpha,\beta \in \WO$ we write $\alpha \ordlneq \beta$ if there is an order-isomorphism from $\alpha$ onto a proper initial segment of $\beta $. We also write $\alpha \ordeq \beta$ if $\alpha $ and $\beta $ are order-isomorphic, and $\alpha \ordleq \beta $ if $\alpha \ordlneq \beta \lor \alpha \ordeq \beta$.
Thanks to the uniqueness of comparison maps, we crucially have

\begin{proposition}
[$\ATR_0$]
    $\ordlneq, \ordleq, \ordeq$ are $\Din 1 1 $.
\end{proposition}
    

% Given $\alpha = (X,\leq) $ and $\beta = (X', \leq')$, a function $f:X\to X'$ is a \emph{homomorphism} $\alpha \to \beta$ if $\forall x,y \in X (x\leq y \limp f(x) \leq' f(y))$.
% We write $\alpha \ordleq \beta$ if there is a homomorphism $\alpha \to \beta$.
% % Note that $\ordleq$ is, a priori, $\Sin 1 1$.\footnote{Optimising the logical complexity of ordinal comparison will not be necessary for this work.}
% We also write $\alpha \ordlneq \beta := \alpha \ordleq \beta \land \lnot \beta \ordleq \alpha$.
% % $\ordlneq$ is, a priori, a Boolean combination of $\Sin 1 1$.


% \begin{remark}
% Our notion of comparison $\ordleq$ is sometimes known as \emph{strong comparison} in the literature, e.g.\ \cite{hirst05:ord-arith-rm-survey}. 
% A weaker notion of ordinal comparison (`weak comparison') requires only that the existence of a homomorphism, i.e.\ an order-preserving map.
% % that the image of a homomorphism is an initial section.
% In the presence of $\ATR$ these are, in fact, equivalent (see \cite{hirst05:ord-arith-rm-survey}), so we shall not concern ourselves further with this difference.
% \end{remark}
\anupam{commented remark above, it seems obsolete and incorrect now (since we follow Simpson's `strong comparison').}

% Given $\alpha = (X,\leq) $ and $\beta = (X', \leq')$, a function $f:X\to X'$ is an \emph{isomorphism} $\alpha\to\beta$ if $\forall x,y\in X (x\leq y \liff f(x)\leq'f(y))$.
% We write $\alpha \ordeq \beta$ if there is an isomorphism $\alpha \to \beta$. 
% Again $\ordeq$ is, a priori, $\Sin 1 1$.

We shall now state a number of well-known facts about comparison, all of which may be found in, e.g., \cite{Simpson99:monograph} or \cite{hirst05:ord-arith-rm-survey}.\footnote{Note that we have intentionally refrained from specifying the `optimal' theories for each statement, for simplicity of exposition}

\begin{proposition}
    [Facts about ordinal comparison, $\PCA$]
    \label{prop:ord-comparison-facts}
    Let $\alpha,\beta,\gamma$ be well-orders.
    We have the following:
    \begin{enumerate}
        \item\label{item:comparison-is-preorder} (Comparison is a preorder)
\begin{enumerate}
    \item $\alpha\ordleq \alpha$
    \item $\alpha \ordleq \beta \ordleq \gamma \limp \alpha \ordleq \gamma$
\end{enumerate}
\item\label{item:comparison-pseudo-antisymmetric} (Comparison is pseudo-antisymmetric) If $\alpha \ordleq \beta \ordleq \alpha $ then $\alpha \ordeq \beta$.
\item\label{item:comparison-total} (Comparison is total) 
$\alpha \ordleq \beta \lor \beta \ordleq \alpha$
\item\label{item:comparison-wf} (Comparison is well-founded)
$\forall F: \Nat\to \WO \,  \exists x\  \lnot F(x+1)\ordlneq F(x)$
    \end{enumerate}
\end{proposition}

% We also have the following standard facts about ordinal arithmetic, e.g.\ \todo{cite other Hirst papers}:

\anupam{commented ordinal arithmetic proposition below.}
% \begin{proposition}
%  [Ordinal arithmetic, $\ATR_0$]
%  \label{prop:ord-arith}
% \begin{enumerate}
%     \item\label{item:succ-ord-arith} (Successor) $ \forall \alpha \in \WO \, \exists \gamma \ordgneq \alpha \, \forall \delta \ordgneq \alpha \  \delta \ordgeq \gamma$
%     \item\label{item:max-ord-arith} (Maximum) $\forall \alpha,\beta \in \WO.\,  \exists \gamma \ordgeq \alpha,\beta. \, \forall \delta \ordgeq \alpha, \beta. \ \delta \ordgeq \gamma$
%     \item\label{item:sup-ord-arith} (Supremum) $\forall F: \Wo \to \WO. \exists \gamma \in \WO . \forall x\  Fx \ordleq \gamma$ \anupam{in fact have uniqueness in $\ATR$, but this is unnecessary}
%     \anupam{need to say more about suprema here}
% \end{enumerate}
% \end{proposition}

% The ordinals $\gamma$ whose existence are asserted in the above items will be referred to as $\succ \alpha, \max(\alpha,\beta), \sup F$ respectively henceforth. \anupam{is this okay?}

We shall assume basic ordinal existence principles, in particular constructions for successor ($\succ \alpha$), addition ($\alpha + \beta$) and maximum ($\max(\alpha,\beta))$, initial segments ($\alpha_b$ for $b\in \alpha$), all definable and satisfying characteristic properties provably in $\ATR_0$.
We shall also make crucial use of a `bounding' principle:
\begin{proposition}
    [$\Sin 1 1 $-Bounding, $\ATR_0$]
    \label{prop:bounding}
    Let $\phi(\alpha) \in \Sin 1 1 $ with $\forall \alpha  (\phi(\alpha)\limp \WO(\alpha))$. 
    Then $\exists \beta \in \WO \forall \alpha (\phi(\alpha) \limp \alpha \ordleq \beta)$.
\end{proposition}
When appealing to Bounding, we use notation such as $\upbnd \phi$ for an ordinal bounding all $\alpha \in \WO$ such that $\phi(\alpha)$.


\anupam{commented positive combinations of operators section below. Now obsolete.}
% \subsection{On positive $\Sin 1 1 $-combinations of $\Sin 1 2$-operators}
% We will need to build up some basic properties of a certain class of $\Sin 1 2 $ formulas which suffice to describe our type structure $\HR$, while at the same time enjoying enough closure properties (in particular under substitution) to allow our exposition to go through.



% An \emph{operator} is an expression $\oper{\vec X}{\vec x} {\phi(\vec X,\vec x)}$. 
% Formally we set $\oper{\vec X}{\vec x} {\phi(\vec X,\vec x)}\, \vec Y \, \vec t$ to be the formula $\phi(\vec Y, \vec t)$.
% As abuse of notation, we shall write such an operator simply as $\phi$, when $\vec X$, $\vec x$ are clear from context.\footnote{Note that the notation is consistent with a form of `$\eta$-conversion' on operator notation.}
% For convenience, we shall use fresh syntactic symbols $\operatorvar, \operatorvari$ etc.\ to vary over operators, e.g.\ to define operator substitution, though we insist that this is purely meta-level notation.

% \begin{remark}
%     [Operators as 3rd order objects]
%     Note that operators are essentially 3rd order objects, taking both set and number arguments. 
%     While we use `operator variables' $\operatorvar, \operatorvari$ at the meta-level, we stop short of formally extending the syntax of arithmetic to a higher type.
%     All this said, the necessity of such notions in the development of fixed point theorems hints at a natural formalisation through \emph{higher-order} reverse mathematics, a subject beyond the scope of this work.
% \end{remark}

% \begin{definition}
% [(Positive) operator combinations]
% Let $\Phi$ be a set of operators.
% $\relSin 1 1 {\Phi}$ is the set of $\Sin 1 1 $ combinations of formulas of the form $ \phi(\vec Y, \vec t)$ for $\phi \in \Phi$.
% % 
% $\posrelS {\Phi}$ is the set of $\relSin 1 1 {\Phi}$ formulas positive in $\Phi$.
% \end{definition}


% In practice we shall simply specify sets of combinators by sets of formulas. E.g.\ we shall identify $\Sin 1 n $ with the set of operators $\oper {\vec X}{\vec x}\phi$ with $\phi \in \Sin 1 1 $.
% In this way we have:
% \begin{proposition}
%     [Closure under substitution, $\PSCA$]
%     \label{prop:posrelSin-closed-under-subs}
%  % If $\phi(\operatorvar) \in \posrelSin 1 1 {\Phi} {\Psi, \operatorvar}$ and $ \psi \in \posrelSin 1 1 { \Phi} {\Psi}$ then $\phi(\psi) \in \posrelSin 1 1 {\Phi} {\Psi}$. %%OLD
% $\posrelS{\Sin 1 2} \subseteq \Sin 1 2$.
% \end{proposition}
% \begin{proof}
% By meta-level induction on $\phi(\vec \operatorvar) \in \posrelS (\vec \operatorvar)$, exploiting positivity and using $\Sin 1 2 $-choice at any alternation of FO and SO quantifiers.
% \end{proof}


% In this way we have natural inclusions:
% \begin{corollary}
%  $ \posrelSin 1 1 {\Sin 1 1 \cup \Pin 1 1 } {\Sin 1 2} = \posrelSin 1 1 {\Sin 1 1 } {\Sin 1 2} \subseteq \Sin 1 2 $.
% \end{corollary}

% \todo{maybe streamline/simplify all of above after checking next subsections}

\subsection{Knaster-Tarski theorem and approximants}
\label{sec:kt-thm-and-approximants-in-psca}
\anupam{commented previous setup based on operators}
% Throughout this section, let $\phi = \phi(X,x)$ be a formula that is equivalent, provably in $\PSCA$, to a $\Sin 1 2$-formula positive in $X$ and a $\Pin 1 2 $ formula positive in $X$.


% Throughout this subsection, let $\phi(X,x) \in \posrelSin 1 1 {\WO, \ordleq } X \subseteq \posrelSin 1 1 {\Sin 1 1} X \subseteq \Sin {1,+} 2 (X)$.
% We shall sometimes suppress the variables $X,x$, writing simply $\phi$ for $\phi(X,x)$, and/or identify $\phi$ with the operator $\oper{ X}{ x} {\phi(\vec X,\vec x)}$ when there is no danger of confusion, simply to lighten the syntax. \anupam{check this later}
% % In fact, all formulas we consider will remain in $\posrelSin 1 1 {\WO, \ordleq} X$.
% We work in $\PSCA$, unless otherwise stated.





% \anupam{all of above is old}

Throughout this section let $\phi(X,x) \in \posnegDin 1 2 {X,\vec X}{\vec Y}$.
We shall typically ignore/suppress the parameters $\vec X, \vec Y$, focussing primarily on $X$ and $x$, and sometimes even write $\phi$ instead of $\phi(X,x)$.
We shall work within $\PSCA$ unless otherwise stated.

% \begin{remark}
%     [Reverse mathematics of fixed point theorems]
%     Let us point out that, while previous work on the reverse mathematics of fixed point theorems exist in the literature, e.g.\ \cite{PengYamazaki17:rev-math-fixed-point-thms-no-kt}, even for the Knaster-Tarski theorem \cite{SatoYamazaki17:rev-math-fixed-point-thms-with-kt}, these results apply to situations when the lattice or space at hand is countable (a subset of $\Nat$).
% Here we require a `non-uniform' version of the theorem where sets themselves are elements of the lattice, under inclusion.
% \end{remark}

\begin{remark}
    [Reverse mathematics of fixed point theorems]
    Let us point out that, while previous work on the reverse mathematics of fixed point theorems exist in the literature, e.g.\ \cite{PengYamazaki17:rev-math-fixed-point-thms-no-kt}, even for the Knaster-Tarski theorem \cite{SatoYamazaki17:rev-math-fixed-point-thms-with-kt}, these results apply to situations when the lattice or space at hand is countable (a subset of $\Nat$).
Here we require a version of the theorem where sets themselves are elements of the lattice, under inclusion.
Our operators are specified \emph{non-uniformly}, namely via positive formulas.
This amounts to a sort of non-uniform `3rd order reverse mathematics' of Knaster-Tarski, peculiar to the powerset lattice on $\Nat$.
\end{remark}




While we can define (bounded) approximants along any well-order simply by appeal to $\ATR$, it will be helpful to retain the well-order (and other free set variables) as a parameter of an explicit formula to be bound later in comprehension instances, and so we give the constructions explicitly here.

% In what follows, let us call any $F\subseteq \Nat \times \alpha$ satisfying 
% \[\forall b \in \alpha \forall y (Fby \limp \exists c<_\alpha b \, \phi(Fc,y))\]
% \emph{$\phi$-recursive} (or \emph{recursive in $\phi$}).

\begin{definition}
[(Bounded) approximants]
If $\WO(\alpha)$ and $a \in \alpha$ we write $\apprx X x {\phi} \alpha (a,x) $ (or even $\phi^\alpha (a,x)$) for:
\[
\exists F\subseteq \alpha \times \Nat \ (\forall b \in \alpha \forall y (Fby  \limp \exists c<_\alpha b\,  \phi (Fc,y)) \land  Fax)
\]
% Notice that this formula is, a priori, $\Sin 1 1$.
We also write,
\[
\begin{array}{r@{\ := \ }l}
     \apprx X x \phi \alpha(x) & \exists a \in \alpha \, \apprx X x\phi \alpha(a,x) \\
     \apprx X x \phi \WO (x) & 
     \exists \alpha ( \WO(\alpha) \land  \apprx X x \phi \alpha(x) ) 
\end{array}
\]
sometimes written simply $\phi^\alpha(x)$ and $\phi^\WO(x)$ respectively.
% Notice that $\phi^\alpha(x)$ is still $\Sin 1 1 $ and $\phi^\WO(x)$ is now (modulo logical equivalence) $\Sin 1 2$.
\end{definition}

% The following result is important:


\begin{proposition}
\label{prop:approx-is-posD12}
If $\WO(\alpha)$ then $\phi^\alpha$ is $\posnegDin 1 2 {\vec X}{\vec Y}$.
In particular, $\phi^\alpha(a,x)$ is equivalent to:
\begin{equation}
    \label{eq:G-phi-inductive-Gax}
\forall G \subseteq \alpha \times \Nat
(\forall b \in \alpha \forall y (\exists c <_\alpha b \phi (Gc,y) \limp Gby) \limp Gax )
\end{equation}
\end{proposition}
\begin{proof}
As written $\phi^\alpha (a,x)$ is a positive $\Sin 1 1 $ combination of $\phi$, so it is certainly $\posnegSin 1 2 {\vec X}{\vec Y}$ by Lemma~\ref{lem:pol-sub}.
So it suffices to prove the `in particular' clause,
since \eqref{eq:G-phi-inductive-Gax} is already a positive $\Pin 1 1 $ combination of $\phi$.% 

For this we note that we can `merge' the two definitions into an instance of $\ATR$, obtaining (after $\phi$-comprehension) some $H$ satisfying for all $b\in \alpha $ and $y$:
\begin{equation}
    \label{eq:H-phi-atr}
    Hby \liff \exists c <_\alpha b\,  \phi(Hc,y))
\end{equation}
Now we show that in fact $Hax \liff \phi^\alpha(a,x) $. For the left-to-right implication, we simply witness the $\exists F$ quantifier in the definition of $\phi^\alpha$ by $H$.
For the right-to-left implication, let $F\subseteq \alpha \times \Nat$ such that:
\begin{equation}
    \label{eq:F-phi-recursive}
\forall b \in \alpha \forall y (Fby \limp \exists c <_\alpha b \phi (Fc,y))
\end{equation}
We show $\forall x(Fax \limp Hax)$ by $\alpha$-induction on $a$:
\[
\begin{array}{r@{\ \implies \ }ll}
     Fax & \exists b<_\alpha a \phi (Fb,x) & \text{by \eqref{eq:F-phi-recursive}} \\
        & \exists b<_{\alpha} a \phi (Hb,x) & \text{by inductive hypothesis} \\
        & Hax & \text{by \eqref{eq:H-phi-atr}}
\end{array}
\]
We may prove $Hax \liff \eqref{eq:G-phi-inductive-Gax} $ similarly, concluding the proof.
\end{proof}



Eventually we will also show that $\phi^\WO \in \posnegDin 1 2 {\vec X}{\vec Y}$ too, but we shall need to prove the fixed point theorem first, in order to appeal to a duality property.
First let us note a consequence of the above proposition:
\begin{corollary}
    [(Bounded) recursion]
    \label{cor:(bdd)-recursion}
   Let $\alpha \in \WO$. We have the following:
    \begin{enumerate}
        \item\label{item:bdd-recursion} (Bounded recursion) $\phi^\alpha(a) = \bigcup\limits_{b<_\alpha a} \phi(\phi^\alpha(b))$.
    I.e.\
    \[\phi^\alpha(a,x) \liff \exists b<_\alpha a \phi(\phi^\alpha(b),x)
    \]
    \item\label{item:recursion} (Recursion) $\phi^\alpha = \bigcup\limits_{\beta \ordlneq \alpha} \phi(\phi^\beta)$, i.e.\ 
    \[
    \phi^\alpha(x) \liff \exists \beta \ordlneq \alpha\, \phi(\phi^\beta, x)
    \]
    \end{enumerate}
\end{corollary}
\begin{proof}
    \ref{item:bdd-recursion} follows immediately from the equivalence between $\phi^\alpha$ and $H$ satisfying \eqref{eq:H-phi-atr} in the preceding proof. 
% 
    For \ref{item:recursion} we first need an intermediate result.
    For $b\in \alpha \in \WO$, let us write $\alpha_b$ for the initial segment of $\alpha$ up to (and including) $b$.
    We show,
    \begin{equation}
        \label{eq:bdd-approximant-equals-approximant-upto-initial-segment}
        \forall x (\phi^\alpha (b,x) \liff \phi^{\alpha_b}(x))
    \end{equation}
    by transfinite induction on $b\in \alpha$.\anupam{discuss induction on arithmetical combinations} 
    We have:
    \[
    \begin{array}{r@{\ \iff \ }ll}
         \phi^\alpha(b,x) & \exists c<_\alpha b\, \phi(\phi^\alpha(c),x) & \text{by \ref{item:bdd-recursion}}\\
         & \exists c<_\alpha b \, \phi (\phi^{\alpha_c},x) & \text{by IH} \\
         & \exists c<_\alpha b \, \phi (\phi^{\alpha_b}(c),x) & \text{by IH} \\
         & \exists b' \in \alpha_b \exists c <_{\alpha_b} b' \phi (\phi^{\alpha_b}(c),x) & \text{set $b'=b$} \\
         & \exists b' \in \alpha_b \phi^{\alpha_b}(b',x) & \text{by \ref{item:bdd-recursion}}\\
         & \phi^{\alpha_b}(x) &\text{by dfn.\ of $\phi^{\alpha_b}$}
    \end{array}
    \]

    Now, turning back to \ref{item:recursion}, we have:
    \[
    \begin{array}[b]{r@{\ \iff \ }ll}
         \phi^\alpha(x) & \exists c \in \alpha \, \phi^\alpha(c,x) & \text{by definition} \\
         & \exists c \in \alpha \exists d<_\alpha c\, \phi(\phi^\alpha(d),x) &\text{by \ref{item:bdd-recursion}} \\
         & \exists c \in \alpha \exists d<_\alpha c \, \phi(\phi^{\alpha_d},x) & \text{by \eqref{eq:bdd-approximant-equals-approximant-upto-initial-segment}} \\
         & \exists c \in \alpha \exists \beta \ordlneq \alpha_c \, \phi(\phi^\beta,x) & \text{by basic ordinal properties \comment{?}} \\
         & \exists \beta \ordlneq \alpha \, \phi(\phi^\beta,x) \tag*{\qed}
    \end{array}
    \]
    \anupam{in penultimate step, might need invariance of approximants under isomorphic ordinals.}
    \renewcommand{\qed}{}
\end{proof}

% Note that, as written, $ \phi^\alpha$ 
% % and $\phi^\WO$ 
% is a positive $\Sin 1 1 $ combination of $\phi$.
% % (with a $\Pin 1 1 $ parameter $\WO$ in the case of $\phi^\WO$).
% Thus $\phi^\alpha$ is $\posnegSin 1 2 {\vec X}{\vec Y}$, provably in $\PSCA$, by Substitution Lemma~\ref{lem:pol-sub}
% and we have access to $\phi^\alpha$
% % and $\phi^\WO$
% as a set under $\ca{\Sin 1 2}$.
% From here we can use the Monotonicity Lemma~\ref{lem:pos-imp-mon} to derive:
\begin{proposition}
    [(Bounded) approximants are inflationary]
\label{prop:approximants-are-inflationary}
% Let $\phi(X,x)$ be arithmetical and positive in $X$. 
Let $\alpha,\beta \in \WO$. We have the following:
\begin{enumerate}
    \item\label{item:bdd-approximants-are-inflationary} $a\leq_\alpha b \limp \forall x (\phi^\alpha(a,x) \limp \phi^\alpha(b,x)) )$
    \item\label{item:approximants-are-inflationary} $\alpha \ordleq \beta \limp \forall x (\phi^\alpha(x) \limp \phi^\beta(x))$
\end{enumerate}
\end{proposition}
\begin{proof}
    \ref{item:bdd-approximants-are-inflationary} follows directly from Bounded Recursion:
    \[
    \begin{array}{rcll}
         \phi^\alpha(a,x) & \implies & \exists c<_\alpha a \, \phi(\phi^\alpha(c),x) & \text{by Corollary~\ref{cor:(bdd)-recursion}.\ref{item:bdd-recursion}}\\
         & \implies & \exists c<_\alpha b\, \phi(\phi^\alpha(c),x) & \text{since $a\leq_\alpha b$}\\
         & \implies &\phi^\alpha(b,x) & \text{by Corollary~\ref{cor:(bdd)-recursion}.\ref{item:bdd-recursion}} 
    \end{array}
    \]
\ref{item:approximants-are-inflationary} follows directly from Recursion:
\[
\begin{array}[b]{r@{\ \implies \ }ll}
     \phi^\alpha(x) & \exists \gamma \ordlneq \alpha \, \phi(\phi^\alpha,x) & \text{by Corollary~\ref{cor:(bdd)-recursion}.\ref{item:recursion}}\\
     & \exists \gamma \ordlneq \beta \, \phi(\phi^\gamma, x) & \text{since $\alpha \ordleq \beta$} \\
     & \phi^\beta(x) & \text{by Corollary~\ref{cor:(bdd)-recursion}.\ref{item:recursion}} 
\end{array}
\]
    \anupam{commented old direct proof below}
    % For \ref{item:approximants-are-inflationary}, let $\alpha\ordleq \beta$ and let $f:\alpha \hookrightarrow \beta$ be an embedding of $\alpha$ onto an initial segment of $\beta$.\anupam{better explain this earlier.}
    % Suppose $\phi^\alpha(x)$, so let $F\subseteq \alpha \times \Nat$ satisfy \eqref{eq:F-phi-recursive} and $a \in \alpha$ such that $Fax$.
    % We define by comprehension $G\subseteq \beta \times \Nat $ by:
    % \begin{equation}
    %     \label{eq:embedding-phi-rec-seq}
    %     Gby\ \liff \ \exists b' \in \alpha \, (fb' = b \land Fb'y)
    % \end{equation}
    % Note that, by injectivity of $f$ \anupam{ref?}, we have that:
    % \begin{equation}
    %     \label{eq:embedding-rec-seq-characterisation}
    %     Fcy \liff G(fc)y
    % \end{equation}
    % We have immediately $G(fa)x$, so it suffices to show,
    % \[
    % \begin{array}{rlll}
    %      & Gby\\
    %      \implies & \exists b' \in \alpha (fb' = b \land Fb'y) & \text{by \eqref{eq:embedding-phi-rec-seq}}\\
    %      \implies & \exists b' \in \alpha (fb'= b \land \exists c'<_\alpha b' \, \phi(Fc',y)) & \text{by \eqref{eq:F-phi-recursive}} \\
    %      \implies & \exists b' \in \alpha (fb'= b \land \exists c'<_\alpha b' \, \phi(G(fc'),y)) &\text{by \eqref{eq:embedding-rec-seq-characterisation}} \\
    %      \implies & \exists c'\in \alpha (fc' <_\beta b \land \phi(G(fc'),y)) \\
    %      \implies & \exists c<_\beta b \, \phi (Gc,y) 
    % \end{array}
    % \]
    % where the final two lines follow from $f$ being an embedding $\alpha \hookrightarrow \beta$. \anupam{actually final line does not seem to rely on this, just inclusion of domains.}
    \renewcommand{\qed}{}
\end{proof}





\anupam{commented bounded approximant collection below}
% \begin{lemma}
% [Bounded approximant collection]
% Let $\psi(X)$ be arithmetical and positive in $X$, and let $\alpha \in \WO$. 
% Then:
% \[
% \psi(\phi^\alpha) \liff \exists b \in |\alpha| 
% \]
% \end{lemma}



\begin{definition}
[Least (pre)fixed points]
% Let $\phi(X,x)$ be an arithmetical formula positive in $X$.
Define $(\mufo X x \phi)(y)$ (or even $(\mu\phi)(y)$) by:
\[
(\mufo X x \phi) (y) := \forall X (\forall x ( \phi(X,x) \limp Xx ) \limp Xy)
\]
\end{definition}
% Note that the above is $\Pin {1,+}1 (\phi)$.
Note that we may treat $\mu\phi$ as a set in $\PSCA$ since $\phi\in \Din 1 2$, and so $(\mu\phi)(y) $ is $ \Pin 1 2$.
By mimicking a text-book proof of the Knaster-Tarski theorem we obtain:
\begin{proposition}
[`Knaster-Tarski']
    \label{prop:kt-for-least-pre-fixed-point}
    $\phi(\mu\phi)= \mu \phi$, i.e.\ 
    \[
    \forall x (\phi(\mu\phi,x) \limp \mu\phi\, x)
    \]
\end{proposition}
\begin{proof}
    % Writing $\phi(X)\subseteq X$ for $\forall x (\phi(X,x) \limp Xx)$, 
    For $\phi(\mu\phi)\subseteq \mu\phi$, note that for any $X$ with $\phi(X)\subseteq X$ we have:
    \[
    \begin{array}{r@{\ \limp \ }ll}
         \mu\phi\, x & Xx & \text{by definition of $\mu\phi$} \\
         \phi(\mu\phi,x) & \phi(X,x) & \text{by Lemma~\ref{lem:pos-imp-mon}}\\ 
         \phi(\mu\phi,x) & Xx & \text{since $\phi(X)\subseteq X$}
    \end{array}
    \]
    Thus $\forall X(\phi(X)\subseteq X \limp \forall x ( \phi(\mu\phi,x) \limp Xx)) $, and so indeed $\phi(\mu\phi)\subseteq \mu\phi$.

    For $\mu\phi \subseteq \phi(\mu\phi)$, we have:
    \[
    \begin{array}[b]{rcll}
         \phi(\mu\phi) &\subseteq & \mu\phi &\text{by above}\\
         \phi(\phi(\mu\phi)) & \subseteq &\phi(\mu\phi) & \text{by Monotonicity Lemma~\ref{lem:pos-imp-mon}}\\
         \phi(\mu\phi) & \subseteq & \mu\phi & \text{by comprehension on $\phi(\mu\phi)$} \tag*{\qed}
    \end{array}
    \]
    \renewcommand{\qed}{}
\end{proof}

\begin{lemma}
    [Closure ordinals]
    \label{lem:closure-ordinals}
    $\exists \alpha \in \WO \, \phi^\WO \subseteq \phi^\alpha$.
    I.e.\ 
    \[
    \forall x (\phi^\WO(x) \limp \phi^\alpha(x))
    \]
\end{lemma}
\begin{proof}
    [Proof sketch]
    We have:\footnote{For the penultimate step, recall that we have access to $\phi^\WO$ as a set, under $\Sin 1 2$-comprehension.}
    \[
    \begin{array}{rll}
         & \forall x (\phi^\WO(x) \limp \exists \alpha \in \WO \phi^\alpha(x)) & \text{by dfn.\ of $\phi^\WO$}\\
         \therefore & \forall x \exists \alpha \in \WO (\phi^\WO(x) \limp \phi^\alpha(x)) & \text{by pure logic}\\
         \therefore & \exists F : \Nat \to \WO \, \forall x (\phi^\WO(x) \limp \phi^{Fx}(x))  & \text{by $\Sin 1 2 $-choice}\\
         \therefore & \forall x (\phi^\WO \limp \phi^{\upbnd F }(x)) & \text{by Prop.~\ref{prop:approximants-are-inflationary}.\ref{item:approximants-are-inflationary}}
    \end{array}
    \]
    where $\upbnd F\in \WO$ is obtained by Bounding,\footnote{To be precise, we apply bounding on the formula $\exists x \forall a (Xa \liff Fxa)$.} Proposition~\ref{prop:bounding}, such that $\forall x \upbnd F \ordgeq Fx$.
\end{proof}

\anupam{commented approximant collection below}
% \begin{lemma}
% [Approximant collection]
% Let $\psi(X)$ be $\Sin 1 1(\vec A)$ and positive in $X$, for $\vec A$ some $\Sin 1 1 \cup \Pin 1 1 $ properties.
% Then:
% \[
% \psi(\apprx Xx \phi\WO) \liff \exists \beta \in \WO \psi (\apprx Xx \phi\beta)
% \]
% \end{lemma}
% \begin{proof}[Proof sketch]
% The right-to-left implication is easy \anupam{any comprehension needed? Idea is functoriality, but requires $\exists $-steps, surely fine.d}, so we focus on the left-to-right implication.
% We proceed by (meta-level) induction on the structure of $\psi(X)$.
% \begin{itemize}
%     \item For $\lor$ and $\land$ cases we simply take the maximum of ordinals obtained by the inductive hypothesis.
%     \item For $\exists$ the same ordinal will do.
%     \item Suppose $\psi(X) = \forall x \psi'(X,x)$. Then we have $\forall x \exists \beta \in \WO \psi'(\apprx X \phi\beta, x)$ by the inductive hypothesis. 
%     By (provably) $\ac{\Sin 1 2}$ we have $\exists F : \Nat \to \WO \forall x \psi'(\apprx X \phi{F(x)},x)$.
%     We can now set $\beta := \sup F$. \oldtodo{check this, might need $F$ to be itself well-ordered} \anupam{actually since $F$ is an $\omega$-sequence, this should be unimportant. Can extract a well-ordered sequence by iterated max, which is not possible for longer sequences.}
% \end{itemize}
% \end{proof}

% \oldtodo{proof}

% \begin{lemma}
% [Successor]
% \oldtodo{need a successor lemma. Could be better to prove expected things about approximants from their bounded variants.}
% \end{lemma}

The main result of this subsection is:
% We can now prove the main LFP characterisation theorem.
\begin{theorem}
[LFP dual characterisation]
\label{thm:muphi-is-phiwo}
% Let $\phi(X,x) $ be a formula positive in $X$. 
$\mu \phi =  \phi^\WO$, i.e.:
\[
% \forall x \left(
% (\mufo X x \phi) (x) \liff \apprx X x \phi\WO(x)
% \right)
\forall x \left(
\mu \phi (x) \liff \phi^\WO(x)
\right)
\]
\end{theorem}
\begin{proof}
For the left-right inclusion, $\mu \phi \subseteq  \phi^\WO$, we show that $\phi^\WO$ is a pre-fixed point of $\oper X x \phi$, i.e.:
\[
\forall x (\phi (\phi^\WO,x) \limp \phi^\WO (x))
\]
By Lemma~\ref{lem:closure-ordinals} let $\alpha $ such that $\phi^\WO = \phi^\alpha$.
We have:
\[
\begin{array}{r@{\ \implies \ }ll}
     \phi(\phi^\WO,x) & \phi(\phi^\alpha,x) & \text{by assumption} \\
     & \phi^{\succ \alpha}(x) & \text{by Cor.~\ref{cor:(bdd)-recursion}.\ref{item:recursion} } \\
     & \phi^\WO (x) & \text{by definition of $\phi^\WO$}
\end{array}
\]
Since we have access to $\phi^\WO$ as a set, under $\ca{\Sin 1 2}$, this indeed yields $\mu \phi \subseteq \phi^\WO$.

For the right-left inclusion, $\phi^\WO \subseteq \mu \phi$, let $\alpha \in \WO$ and we show:
\[
\forall a \in \alpha ( \phi^\alpha(a,x) \limp \mu \phi\, x)
\]
by transfinite induction on $a \in \alpha$.
For logical complexity, recall that we indeed have access to $\mu\phi$ as a set, since $\phi\in \Din 1 2$ so $\mu\phi \in \Pin 1 2$.
We have:
\[
\begin{array}[b]{r@{\ \implies \ }ll}
     \phi^\alpha(a,x) & \exists b<_\alpha a \, \phi(\phi^\alpha(b),x) & \text{by Corollary~\ref{cor:(bdd)-recursion}.\ref{item:bdd-recursion}}\\
     & \exists b<_\alpha a \, \phi(\mu\phi,x) & \text{by inductive hypothesis} \\
     & \phi(\mu\phi,x) & \text{by vacuous quantification} \tag*{\qed}
\end{array}
\]
\renewcommand{\qed}{}
\end{proof}




One of the main consequences of the above result is:
\begin{corollary}
    [$\PSCA$]
    \label{cor:lfp-of-posD12-is-posD12}
    $\mufo X x \phi$ is $\posnegDin 1 2 {\vec X}{\vec Y}$.
\end{corollary}

\anupam{commented old nu stuff below. obsolete.}
% \subsection{Using duality to optimise complexity}


% \anupam{So, actually, we really do need that $\mufo X x \phi$ is $\Din 1 2 $ in $\PSCA$. need a little lemma for this, but goes through via pre/post definition and main characterisation above.}

% The final part of this section is to establish a suitable logical complexity bound on fixed points. 
% For this it is useful to reconsider the dual of a least fixed point, namely the greatest fixed point, within arithmetic (now in the classical setting).

% Given a formula $\phi(X,x)$ let us write $(\nufo X x \phi)(a) $ (or even $(\nu \phi) (a)$) for:
% \[
% \exists X (\forall x (Xx \limp \phi(X,x)) \land Xa)
% \]

% Let us furthermore write $\bar \phi(X,x)$ for $\lnot \phi(\lnot X,x)$. 
% We have the following from almost purely logical rewriting:
% \begin{proposition}
%     [$\ACA$]
%     $(\mufo X x \phi)(a) \liff \lnot (\nufo X x {\bar \phi})(a)$.
% \end{proposition}

% Now, we can repeat almost all of this section's development to characterise greatest fixed points in terms of the \emph{intersection} of co-approximants, again provably $\Sin 1 2$. \anupam{say more}
% By the characterisation theorem, we thus obtain a powerful consequence on logical complexity:

% \begin{corollary}
%     [$\PSCA$]
%     $\mufo X x \phi$ is $\posnegDin 1 2 {\vec X}{\vec Y}$.
% \end{corollary}

Let us point out that the above result amounts to a partial arithmetisation of purely descriptive characterisation of $\mu$-definable sets in Lubarsky's work \cite{Lubarsky93}.
Note also that, while we (in particular) obtain a $\Delta^1_2$ bound on fixed point formulas, we crucially required $\CA{\Pi^1_2}$ to prove this, consistent with results of \cite{Moellerfeld02:phd-thesis} that we shall exploit in the next section to fill in our `grand tour'. 

\subsection{Arithmetising the totality argument}




Thanks to the results of this section, we may duly formalise the type structure $\HR$ from Section~\ref{sec:totality} within $\PSCA$ and prove basic properties.
% 
% 
Throughout this section we shall identify terms with their codes. In the case of $\clomuLJneg$ and subsystems, we note that terms may still be specified finitely thanks to regularity of coderivations in $\cmuLJneg$, and that syntactic equality is verifiable already in $\RCA$ between different representations~\cite{Das21:CT-preprint,Das21:CT-fscd}.
Let us also note that the various notions of reduction, in particular $\convertsnorec$ and $\convertsnoreceta $ are $\Sin 0 1$ relations on terms. (Note here that it is convenient that $\convertsnoreceta$ is formulated using only the extensionality rule, for $\Sin 0 1$, rather than being fully extensional.)

Let us recall the $\HR$ structure from Section~\ref{sec:totality}. 
Note that, as written in Definition~\ref{def:hr},
these sets, a priori, climb up the analytical hierarchy due to alternation of $\mu$ and $\arrow$.
This is where the results of the previous section come into play and serve to adequately control the logical complexity of $\HR$.

% Thus induction principles on them (such as those required for the totality argument for circular proofs) is not, a priori, available to us within $\PSCA$.\anupam{make this remark rather the motivation for the previous section's development}

% However, we can duly represent this type structure by referring only to explicit ordinal approximants of the fixed points, using some of the technology we developed in the previous section.
% Let us make this definition explicitly here:


% \todo{rewrite introductory text above later}

\begin{definition}
[Type structure, formalised]
We define the following second-order formulas:
\begin{itemize}
    % \item $\hr \emptyset (t) := \bot$
    % \item $\hr A (t) := A(t)$\oldtodo{don't need this now that we treat fixed points differently?}
    % \item $\hr \unit (t) := t \betae \identity $\footnote{Recall that $\betae$ is $\Sin 0 1 $ so, in particular, arithmetical.}
    \item $\hr X (t) \dfn Xt$
    \item $\hr \nat (t) \dfn  \exists n\,  (t \convertsnoreceta \numeral n) $ 
    \item $\hr {\sigma \to \tau} (t) \dfn \forall s ( \hr\sigma (s) \limp \hr \tau (ts))$
    \item $\hr{\sigma \times \tau} (t) \dfn \hr \sigma (\proj 0 t) \land \hr \tau (\proj 1 t) $
    % \item $\hr{\sigma + \tau} (t) := (\project 0 t \betae \enc\false \land \hr\sigma(\project 1 t)) \lor (\project 0 t \betae \enc\true \land \hr\tau(\project 1 t))$
    % \item $\hr{\sigma^\alpha(\emptyset)}(t):= \hr{\sigma}^\alpha (t)$ (this is optional)
    % \item $\hr{\mu X \sigma}(t) \dfn \apprx X x{\hr{\sigma}(x)}\WO(t)$ 
    \item  $\hr {\mu X \sigma} (t) \dfn (\mufo X x {\hr \sigma (x)}) (t)$
\end{itemize}
\end{definition}

Weak theories such as $\RCA$ are able to verify closure of $\clocmuLJneg$ under conversion by a syntax analysis (see~\cite{Das21:CT-fscd}).
That each $\hr\sigma$ itself is closed under conversion, cf.~\ref{prop:hr-closed-under-conversion}, is also available already in $\RCA$ by (meta-level) induction on the structure of $\sigma$.
More importantly, and critically, we have as a consequence of the previous section:

\begin{corollary}
    \label{cor:hr-is-D12}
    Let $\sigma$ be positive in $\vec X$ and negative in $\vec Y$.
    $\hr \sigma$ is $\PSCA$-provably $\posnegDin 1 2 {\vec X}{\vec Y}$. 
    % equivalent to both:
    % \begin{enumerate}
    %     \item a $\Sin 1 2$ formula positive in $\vec X$ and negative in $\vec Y$; and,
    %     \item a $\Pin 1 2 $ formula positive in $\vec X$ and negative in $\vec Y$.
    % \end{enumerate}
\end{corollary}
\begin{proof}[Proof sketch]
    We proceed by induction on the structure of $\sigma$, for which the critical case is when $\sigma$ has the form $\mu X \tau$.
    By definition $\tau$ is positive in $X$, and so by IH $\hr{\tau} \in \posnegDin 1 2 {X,\vec X}{\vec Y}$.
    We conclude by Corollary~\ref{cor:lfp-of-posD12-is-posD12}.
\end{proof}


% In particular, by induction on type we have:

% \begin{corollary}
%     \label{cor:hr-is-D12}
%     Let $\sigma$ be positive in $\vec X$ and negative in $\vec Y$.
%     $\hr \sigma$ is $\PSCA$-provably $\posnegDin 1 2 {\vec X}{\vec Y}$. 
%     % equivalent to both:
%     % \begin{enumerate}
%     %     \item a $\Sin 1 2$ formula positive in $\vec X$ and negative in $\vec Y$; and,
%     %     \item a $\Pin 1 2 $ formula positive in $\vec X$ and negative in $\vec Y$.
%     % \end{enumerate}
% \end{corollary}
% \begin{IEEEproof}[Proof sketch]
%     We proceed by induction on the structure of $\sigma$, for which the critical case is when $\sigma$ has the form $\mu X \tau$.
%     By definition $\tau$ is positive in $X$, and so by IH $\hr{\tau} \in \posnegDin 1 2 {X,\vec X}{\vec Y}$.
%     We conclude by Corollary~\ref{cor:lfp-of-posD12-is-posD12}.
% \end{IEEEproof}

Now, recall that a function $f:\Nat^k \to \Nat$ is \emph{provably recursive} in a theory $T\subseteq \SOPA$ if there is a $\Sin 0 1 $-formula $\phi_f(\vec x,y)$ with:
\begin{itemize}
    \item $ \nmod \models \phi_f (\vec {\numeral m} , \numeral n) $ if and only if $f(\vec m) = n$; and,
    \item $T \proves \forall \vec x \exists y \phi_f (\vec x, y)$.
\end{itemize}

We may formalise the entire totality argument within $\PSCA$, in a similar fashion to analogous arguments in \cite{Das21:CT-fscd,KuperbergPP21systemT,Das21:CT-preprint}, only peculiarised to the current setting.
This requires some bespoke arithmetical arguments and properties of our type structure.
% , explained further in Appendix~\ref{sec:formalised-totality}.
As a consequence we shall obtain:

\begin{theorem}
\label{thm:cmuLJneg-to-psca}
Any $\cmuLJneg$-representable function on natural numbers is provably recursive in $\PSCA$.
\end{theorem}


We give some more details for establishing the above theorem. In particular, we apply the fixed point theorems within second-order arithmetic from Section~\ref{sec:kt-thm-and-approximants-in-psca} to arithmetise the totality argument in Section~\ref{sec:totality} within $\PSCA$.
We shall focus on explaining at a high level the important aspects of the formalisation, broadly following the structure of Section~\ref{sec:totality}.




Notice that it will not actually be possible to prove the Interpretation Theorem~\ref{thm:prog-implies-hr} uniformly, due to G\"odelian issues.
Instead we will demonstrate a non-uniform version of it:
\begin{theorem}
    [Nonuniform Interpretation, formalised]
    \label{thm:interpretation-nonuniform-formalised}
    Let $\der:\Sigma \seqar \tau$.  
    Then
    $\PSCA \proves \forall \vec s \in \hr\Sigma .\  \der \vec s \in \hr \tau $.
\end{theorem}


Note that, from Theorem~\ref{thm:interpretation-nonuniform-formalised} above, Theorem~\ref{thm:cmuLJneg-to-psca} will follow directly via a version of Corollary~\ref{cor:cmuLJ-represents-only-total-fns} internal to $\PSCA$.
In particular, when $\Sigma$ contains only $\N$s and $\tau$ is $\N$, note that the statement of Theorem~\ref{thm:interpretation-nonuniform-formalised}, $\forall \vec s \in \hr\Sigma .\  \der \vec s \in \hr \tau $, is indeed $\Pin 0 2$.

The rest of this section is devoted to justifying \Cref{thm:interpretation-nonuniform-formalised} above.
Let us fix $P:\Sigma \seqar \tau$ for the remainder of this section.


\subsubsection*{Formalising monotonicity and transfinite types}
All the technology built up in Subsection~\ref{subsec:monotonicity-and-transfinite-types} has been appropriately formalised in the previous Section~\ref{sec:kt-thm-and-approximants-in-psca}.

\subsubsection*{Formalising closures and priorities}
All the notions about closures and priorities in the totality argument involve only finitary combinatorics and are readily formalised within $\RCA$.

\subsubsection*{Formalising ordinal assignments}
We define ordinal assignments within $\PSCA$ relative to some $\vec \alpha$ varying over $\WO$. 
% 
% to obtain a form of `approximant collection':
% \[
% \forall x \phi^\WO(x) \limp \exists \alpha \in \WO \forall x \phi^\alpha(x)
% \]
% Ordinal assignments are easily defined by simply witnessing existential quantifiers in $\phi^\WO$ by appropriate well-orders.
What is important is to establish the `positive and negative approximants' Proposition~\ref{prop:pos-neg-approximants}. 
However this follows directly from the Closure Ordinal Lemma~\ref{lem:closure-ordinals} and the Monotonicity Lemma~\ref{lem:pos-imp-mon}.
% However this follows directly from approximant collection and well-foundedness of well-order-comparison.
% \oldtodo{important here that we have $\notin$ for the negative assignments, so that Approximant Collection is duly applicable.}

\subsubsection*{Formalising reflection of non-totality.} 

Since $P$ has only finitely many distinct lines, and so only finitely many distinct formulas, we will need to consider only finitely many distinct $\hr\sigma$ henceforth which we combine to describe to establish the reflection of non-totality Definition~\ref{def:refl-non-totality} all at once within SO arithmetic.
% In particular using ordinal existence principles, we can take some master ordinal that exceeds the closure ordinals, cf.~\ref{lem:closure-ordinals}, of all 
Working in $\PSCA$, let us point out that the description of $P', \Sigma', \tau',\vec s'$ from $\rrule$ and $P$'s (finitely many) sub-coderivations $\vec P$ is recursive in the following oracles:
\begin{itemize}
    \item $\hr\sigma$, for each $\sigma $ occurring in $P$, which is $\Din 1 2 $ by Corollary~\ref{cor:hr-is-D12}; and,
    \item finding the least $\alpha \in \WO$ such that $ \hr{\sigma^\alpha}(t)$; this is $\Din 1 2 $ by fixing the closure ordinal, cf.~\ref{lem:closure-ordinals}, say $\gamma \in \WO$ of $\mu \hr\sigma$, and searching for the least appropriate $a\in \gamma$ instead; again this is $\Din 1 2$;
    \item finding an inhabitant $s$ of some nonempty $\hr\sigma$, for which we can simply take the `least' (seen as a natural number coding it), and so is $\Din 1 2$.
\end{itemize}

Consequently we have that Definition~\ref{def:refl-non-totality}, the description of $P', \Sigma', \tau',\vec s'$ from $\rrule, \vec P$, is indeed a $\PSCA$-provably $\Din 1 2$ formula in $P', \Sigma', \tau',\vec s',\rrule, \vec P$.

\anupam{point out we need KT theorem for $\mu$-right?}

\subsubsection*{Formalising the non-total branch construction}

The `non-total' branch constructed in the proof of the Interpretation Theorem~\ref{thm:prog-implies-hr} is recursive in the `reflecting non-totality' definition, and so we have access to it as a set within $\PSCA$.


One subtlety at this point is that $\PSCA$ needs to `know' that $\der$ is indeed progressing.
However earlier work on the cyclic proof theory of arithmetic~\cite{Simpson17:cyc-arith,Das20:ca-log-comp}, building on the reverse mathematics of $\omega$-automaton theory \cite{KMPS19:buchi-rev-math}, means that this poses us no problem at all:

\begin{proposition}
    [Formalised cyclic proof checking~\cite{Das20:ca-log-comp}]
    $\RCA$ proves that $P$ is progressing, i.e.\ that every infinite branch has a progressing thread.
\end{proposition}

From here we readily have that there is a progressing thread $(\rho_i)_{i\geq k}$ and inputs $r_i$ along along the non-total branch.
The assignment of ordinals $(\vec\alpha_i)_{i\geq k}$ along the branch follows by the `positive and negative approximants' result (formalisation of \Cref{prop:pos-neg-approximants}).
The verification that $\vec \alpha_{i+1}\leq \vec \alpha_i$ for non-critical steps follows by inspection of the `reflecting non-totality' definition (formalisation of \Cref{def:refl-non-totality}).
Finally, for the critical fixed point unfolding, we need that $\alpha_{in}$ is a successor. 
For this we need that the approximant at a limit ordinal is a union of smaller approximants, for which we use Recursion, cf.~Corollary~\ref{cor:(bdd)-recursion}.

\medskip

This concludes the argument for Theorem~\ref{thm:interpretation-nonuniform-formalised}, and so also of Theorem~\ref{thm:cmuLJneg-to-psca}.